
\illustration{0.5}{colours}{Image taken from \href{https://commons.wikimedia.org/wiki/File:Wordle_Emoji_Screenshot.png}{commons.wikimedia.org}.}

Ingfríður is testing her new pet project website, Passwordle.
The rules should be rather familiar for those who have played
plenty of wordle, but let us review them here. The website picks
a secret password that the user then has to guess. The password
will be $N$ characters in length and the user will guess until
they get it right, getting a higher score for fewer guesses.
Each guess must be a string of $N$ characters. For each character
in the guess, it will be coloured one of three colours. If that
character matches the password character in the same position,
the colour is green. If the character does not match, but that
character is in the password somewhere else, the colour is
yellow. Otherwise it is gray.

Ingfríður is of course very good at her own game, so she will
always guess a password that could be the hidden password.
That is to say her guess will always match the previous clues.
Furthermore she knows her program never generates passwords with
repeated characters, since that's insecure, and will incorporate
this knowledge into her guesses.

You now receive some screenshots from her testing progress, but
they got so terribly compressed you can only make out the colours
and not the text itself. Furthermore it seems that she didn't 
manage to make any progress at all. She didn't get a single green
square in the entire game and never found any more characters than
she did in her very first guess, and quit
out of frustration. So the number of yellow squares is the same
on each row. Given this info, can you reconstruct a
sequence of guesses she could have made? Or is it impossible and
her program must be wrong? You may assume that Ingfríður never
makes any mistake when playing, only when programming.

